\section{The Five Primitives}
\label{sec:primitives}
For each primitive we give its \emph{mechanism}, its \emph{negative definition} (what it never does, and which primitive does), and its \emph{design space}. The negative definitions make the five mutually exclusive: each piece of a harness has one home, assigned by its function rather than by its name or its location in the code.

\subsection{\prim{Input}: capture what comes in}
\textbf{Mechanism.} Capture what enters the agent and hand it to \prim{Context}. There are two sources: a person (a message, a file, an interrupt, an answer to a question the model asked) and the world (the result of an action, an event, a scheduled trigger).
\textbf{Never.} \prim{Input} never decides what the model is shown; that is \prim{Context}. It never runs anything; that is \prim{Output}. An answer a person gives that the \emph{harness} consumes, such as an approval, is a \prim{Control flow} decision; an answer the \emph{model} consumes, returned as a tool result, is input.
\textbf{Design space.} The source (terminal, chat channel, webhook) and the kind (text, files, results, interrupts).

\subsection{\prim{Context}: assemble what the model is shown}
\textbf{Mechanism.} Assemble, for one call, everything the model is shown: instructions, the tools' descriptions, the conversation so far, and anything retrieved from memory or the environment; and fit it to the window by truncation, compaction or selection. \prim{Context} is the agent's working memory, rebuilt on every call.
\textbf{Never.} \prim{Context} never calls the model; that is \prim{Model interface}. It never acts on the world; that is \prim{Output}. Reading a memory store is \prim{Context}; writing one, if the model requests it, is an action and so \prim{Output}. Code that decides what the model sees is \prim{Context} even when it sits inside a tool.
\textbf{Design space.} What the instructions hold; how tools are described; how long-term memory comes in (preloaded, retrieved, or searched by the agent); the roles each piece is placed in; what happens when the window fills.

\subsection{\prim{Model interface}: reach the model}
\textbf{Mechanism.} Send the assembled context to the model and bring back its response in a form the rest of the harness can read: words, and structured requests to act. This includes the client, the wire format, the translation to the model's chat template, streaming, retries, and switching to a backup model when a request could not be served.
\textbf{Never.} \prim{Model interface} never decides what the request contains; that is \prim{Context}. It never acts on the response; that is \prim{Output}. It receives a stream; rendering each fragment is \prim{Output}. A change of model because the request could not be served is \prim{Model interface}; a change of model keyed to the work (planning, a cheap model for a side task) is \prim{Control flow}.
\textbf{Design space.} A hosted API or local inference; one provider or many; native tool-call formats or calls parsed from text; retry policy and backup models. \prim{Model interface} is the only primitive that knows which model is inside, which is why changing the model means changing it alone.

\subsection{\prim{Output}: carry the response out}
\textbf{Mechanism.} Handle what the model returned: show its words to a person, and carry out each requested action, collecting its result. The tool set, the executor, and the environment an action runs in are \prim{Output}.
\textbf{Never.} \prim{Output} never decides whether to go again; that is \prim{Control flow}. It never assembles the next request; that is \prim{Context}. The model only ever \emph{requests} an action; \prim{Output} decides whether and how it happens, which is where sandboxes and never-lists live. Several tool calls from one response run concurrently: \prim{Output}. Several \emph{model} calls run concurrently: \prim{Control flow}.
\textbf{Design space.} The tool set, from one shell to dozens of specific tools; the action format, from structured calls to model-written code \cite{wang2024codeact}; the execution environment; and the presentation of results, which \cite{yang2024sweagent} show affects performance.

\subsection{\prim{Control flow}: decide what happens next}
\textbf{Mechanism.} Define the shape of the loop: send results back and go again, hand back to the person, or stop. Step and token budgets, approval gates the harness consumes, delegation to sub-agents, and routing to a model keyed to the work are all \prim{Control flow}.
\textbf{Never.} \prim{Control flow} never captures input, assembles context, calls the model or runs a tool. It sequences the other four. A workflow in which code fixes the order of model calls \cite{anthropic2024effective} and an agent in which the model directs its own process differ here and only here.
\textbf{Design space.} An open loop, a fixed pipeline or an explicit graph; sub-agents; approval modes; stopping rules.

\paragraph{Necessity} Each primitive answers a question no other answers. The five minimal harnesses published with the paper were built one primitive at a time around a fixed model and task, and each version fails in a specific way without the next: without \prim{Output} the model cannot act, without \prim{Control flow} a result goes nowhere, without \prim{Context} the model spends actions finding out where it is. This shows that each primitive removes a failure the others do not, not that no sixth could be needed; \S\ref{sec:deconstruction} tests that.
